\section{Tweedie's Formula}

\subsection{公式陈述与直觉}

Tweedie's formula 是理解 score function 与扩散模型联系的关键数学工具。它回答了这样一个问题：如果我们知道边缘分布 $p(x)$ 和似然分布 $p(x \mid \mu)$，能否从观测 $x$ 推断出隐变量 $\mu$ 的期望？

\begin{importantbox}{Tweedie's Formula}
考虑以下生成过程：先从先验分布 $p(\mu)$ 采样均值 $\mu$，然后从似然分布 $p(x \mid \mu) = \mathcal{N}(x;\; \mu, \sigma^2 I)$ 采样观测 $x$。令边缘分布为 $p(x) = \int p(x \mid \mu)\, p(\mu)\, d\mu$，则后验均值为：
$$\mathbb{E}[\mu \mid x] = x + \sigma^2 \nabla_x \log p(x)$$
\end{importantbox}

\begin{knowledgebox}{Tweedie's Formula 的直觉}
这个公式有非常优美的直觉解释：
\begin{itemize}
    \item 如果我们只知道观测 $x$ 而没有其他信息，对 $\mu$ 的最朴素估计就是 $x$ 本身。
    \item 但 score function $\nabla_x \log p(x)$ 提供了一个``修正方向''：它指向高概率密度区域。
    \item Tweedie's formula 告诉我们，最优的后验均值估计就是在朴素估计 $x$ 的基础上，沿着 score 方向做一步校正，校正幅度与噪声方差 $\sigma^2$ 成正比。
\end{itemize}
\end{knowledgebox}

\subsection{简单高斯案例验证}

讲者用一个简单的例子来验证 Tweedie's formula 的正确性和直觉：

设 $p(\mu) = \mathcal{N}(0, 1)$，$p(x \mid \mu) = \mathcal{N}(x;\; \mu, \sigma^2)$。由高斯分布的性质（先验和似然都是高斯时，边缘分布也是高斯），可得：

$$p(x) = \mathcal{N}(0, 1 + \sigma^2)$$

因此 score function 为：

$$\nabla_x \log p(x) = -\frac{x}{1 + \sigma^2}$$

代入 Tweedie's formula：

$$\mathbb{E}[\mu \mid x] = x + \sigma^2 \cdot \left(-\frac{x}{1+\sigma^2}\right) = x \cdot \frac{1}{1+\sigma^2}$$

\begin{knowledgebox}{直觉验证：收缩估计}
结果 $\mathbb{E}[\mu \mid x] = \frac{x}{1+\sigma^2}$ 具有``收缩''（shrinkage）效应：
\begin{itemize}
    \item 当 $x > 0$ 时，$\mathbb{E}[\mu \mid x] < x$：估计的均值比观测值更靠近零（先验均值的方向）。
    \item 当 $x < 0$ 时，$\mathbb{E}[\mu \mid x] > x$：同理，向零的方向收缩。
    \item $\sigma^2$ 越大（噪声越大），收缩越强：在高噪声下，我们更信任先验。
    \item $\sigma^2 \to 0$ 时，$\mathbb{E}[\mu \mid x] \to x$：噪声消失时，观测本身就是最好的估计。
\end{itemize}
这完全符合贝叶斯推断的直觉。
\end{knowledgebox}

\subsection{在扩散模型中的应用}

将 Tweedie's formula 应用到扩散模型的前向跳跃分布 $q(x_t \mid x_0) = \mathcal{N}(\sqrt{\bar{\alpha}_t}\, x_0, (1-\bar{\alpha}_t)I)$，其中``均值''角色由 $\sqrt{\bar{\alpha}_t}\, x_0$ 扮演，``方差''为 $\sigma_t^2 = 1-\bar{\alpha}_t$，可以得到：

$$\mathbb{E}[\sqrt{\bar{\alpha}_t}\, x_0 \mid x_t] = x_t + (1-\bar{\alpha}_t)\, \nabla_{x_t} \log q(x_t)$$

整理后：

\begin{importantbox}{Tweedie's Formula 在扩散模型中的形式}
$$\mathbb{E}[x_0 \mid x_t] = \frac{1}{\sqrt{\bar{\alpha}_t}} \left(x_t + (1-\bar{\alpha}_t)\, \nabla_{x_t} \log q(x_t)\right)$$
这表明：知道边缘 score function $\nabla_{x_t} \log q(x_t)$ 就等价于知道在给定 $x_t$ 下 $x_0$ 的最优估计。Score function 包含了``去噪''所需的全部信息。
\end{importantbox}

\begin{warningbox}{条件 Score 与边缘 Score 的区别}
这里出现的 $\nabla_{x_t} \log q(x_t)$ 是\textbf{边缘}分布的 score，而非条件分布 $q(x_t \mid x_0)$ 的 score。条件 score 可以直接用高斯分布公式计算，但边缘 score 需要对所有可能的 $x_0$ 做积分平均。在训练时，神经网络通过大量样本的平均来隐式学习边缘 score。
\end{warningbox}

\subsection{本章小结}

Tweedie's formula 在``观测即均值加噪声''的模型中，给出了后验均值的优雅闭式解：只需在观测值上加上 score function 方向的校正。在扩散模型中，它将 score function 与去噪操作（从 $x_t$ 恢复 $x_0$ 的期望）紧密联系起来，提供了理解 score-based 方法的关键桥梁。

